% ===================== ABOUT THE AUTHOR =====================
\clearpage
\thispagestyle{plain}
\phantomsection
\addcontentsline{toc}{chapter}{About the Author}
\markboth{About the Author}{About the Author}

\begin{center}
{\footnotesize\color{ufogreenink}\scshape About the Author}\par\vspace{3pt}
{\dispfont\LARGE\color{inkblue} EVAN SVENDSEN}\par\vspace{5pt}
\end{center}
\greenrule
\vspace{13pt}\par

\begin{center}
\adjustimage{max size={4.1in}{3.2in}}{author.jpg}
\end{center}
\vspace{3pt}\par
\noindent\begin{minipage}{\linewidth}
{\RaggedRight\scriptsize\color{steel} The author at the Area 51 exhibit --- \emph{Myth or Reality}. Thumbs up
for the myth; the reality is in your hands.\par}
\end{minipage}\par\vspace{13pt}

\noindent\textbf{Evan Svendsen} is a game designer and software developer from Didsbury, Alberta.
He spent 3 years at the University of Lethbridge studying computer science before later attending
Lethbridge Polytechnic years later to complete a certificate in Virtual \& Augmented Reality
Design.. He works mostly in C\# and Unity, and has spent more than a
decade building and publishing independent games through Tech Skull Studios --- the small
Lethbridge outfit he helped start in 2013, whose titles have shipped on Android and iOS. His
six-year solo project \emph{Beat The Markets} --- a turn-based digital board game about investing,
built for two to six players --- is on Steam at
\srclink{https://store.steampowered.com/app/2095440/Beat_The_Markets/}. He also
writes on technology, transit and the strange at \srclink{https://evansvendsen.substack.com}.

\vspace{4pt}\par
\figwrap[r]{0.28}{The book that started it. Don L. Wulffson's \emph{Aliens: Terrifying
Extraterrestrial Tales}, a thin paperback of retold abduction and encounter stories aimed at
children, was the author's first contact with the subject at eight years old. Nothing in it is
evidence of anything. It is, however, an honest record of how the material reaches most people:
as a story, told well, long before anybody mentions that it is supposed to be true.}{}{fig:wulffson}%
\noindent None of which explains why he was reading about this in the first place, and the honest
answer is a library paperback. He was eight years old when he found Don L. Wulffson's
\emph{Aliens: Terrifying Extraterrestrial Tales} --- a slim collection of encounter and abduction
stories rewritten for children, sold on the promise that it was all real. He read it until the
spine gave out. It did not make him a believer and it did not make him a sceptic; what it did was
put the question in his hands at an age when nobody had yet told him which side of it he was
supposed to be on. Everything in the preceding twenty chapters --- the archives, the statistics, the
radar returns, the long unglamorous work of finding out which stories survive being checked ---
began as an attempt to find out whether that paperback had been telling him the truth. Most of it
had not been. That turned out to be a far more interesting result than a yes would have
been, and it is the reason this book exists at all: not to debunk the eight-year-old, but to
finally give him the sources.

\vspace{4pt}\par
\Needspace*{3.9in}
\begin{center}
\adjustimage{max size={3.8in}{2.9in}}{author_grad.jpg}
\end{center}
\vspace{3pt}\par
\noindent\begin{minipage}{\linewidth}
{\RaggedRight\scriptsize\color{steel} Convocation day. A credential is a reason to check someone's
working, not a reason to skip it --- a standard the author holds himself to as readily as he holds
it to everyone quoted in Chapter~\ref{ch:intro}.\par}
\end{minipage}\par\vspace{13pt}

\vspace{4pt}\par
\noindent The training is on the record as well as in the work. He returned to Lethbridge
Polytechnic as an adult student and finished with a 4.0 term grade point average, named to the
President's List --- Honours with Great Distinction --- for both the Fall 2024 and Winter 2025
terms. He mentions it here for one reason only, and it is the same reason the sources in this book
are printed in full rather than summarised: a reader is entitled to know whether the person asking
them to check arithmetic, read a chi-square table and follow an argument about instrument error has
ever been made to do any of it under supervision.

\vspace{6pt}\par
\Needspace*{4.6in}
\begin{center}
\adjustimage{width=\textwidth}{plist_f24.jpg}\par\vspace{6pt}
\adjustimage{width=\textwidth}{plist_w25.jpg}
\end{center}
\vspace{3pt}\par
\noindent\begin{minipage}{\linewidth}
{\RaggedRight\scriptsize\color{steel} President's List, Lethbridge Polytechnic --- Honours with
Great Distinction, Fall 2024 (awarded 20 December 2024) and Winter 2025 (awarded 5 May 2025), term
grade point average 4.0 in each. Author's own records.\par}
\end{minipage}\par\vspace{13pt}

\noindent Game development turned out to be unexpectedly good training for this book. A designer spends his
working life asking what a system is actually doing rather than what it appears to be doing,
reproducing a bug before believing it exists, and separating the interesting explanation from the
correct one. Ufology, as the preceding twenty chapters have argued at length, suffers badly from a
shortage of exactly that habit.

He flies, too, which matters more to this book than it might sound. He holds a pilot's licence with
Transport Canada and spent six years in the Royal Canadian Air Cadets. A licence is
not an argument, and it does not settle anything on its own --- but a great deal of ufology's
best-loved testimony is aviation testimony, and reading it is easier when you already know what an
aircraft does, what a compass does when it is unhappy, and how thoroughly a pilot in cloud can lose
track of where he is and how long he has been there. Section~\ref{sec:pilots} and the Bermuda
Triangle of Section~\ref{sec:bermudatriangle} are both, in the end, questions about instruments and
the people reading them. Knowing that spatial disorientation is a normal hazard rather than an
insult to the person experiencing it is why those sections are sceptical without being
unkind.

He also owes one of this book's arguments to a hobby, and it would be dishonest to leave it out.
The author is a furry --- a member of the anthropomorphic-animal fandom, a community that has spent
decades doing nothing but designing, drawing, wearing and arguing about characters with animal heads
and human bodies. Twenty years of that does something to the eye. When he first started reading the
ancient-astronaut literature on Egyptian iconography, the claim that a jackal head on a human frame
must record a real hybrid creature did not read as startling; it read as a category error he had
watched people make about his own hobby for years. A fandom that produces tens of thousands of
animal-headed humanoids a year, none of them memories of anything, is not a control group for
ancient Egypt, and it should not be dressed up as one --- a self-aware modern art culture tells you
nothing about what a Nineteenth Dynasty priest thought he was doing with a jackal. What it does do
is dispose of one premise. If hybrid imagery can be mass-produced by people who know perfectly well
that no such animal exists, then a hybrid picture cannot on its own be evidence of a hybrid
creature. That is a narrow result and it deserves to be kept narrow. Its whole use is that it makes
the actual question visible: not \emph{why does this figure have a jackal's head} but \emph{what is
the picture for}, which is the question Section~\ref{sec:animalheads} and
Section~\ref{sec:animalpeople} spend their pages answering. The hobby did not supply the evidence
--- the captions, the mask finds and the absent skeletons did that --- but it did supply the
instinct to distrust the literal reading, and that instinct arrived a good deal earlier than it
would have otherwise. Insight is allowed to come from unserious places. It only has to survive the
checking afterwards. He notes, with some amusement, that the machines have since taken the hobby's
side: the image-generators of Section~\ref{sec:ai} now mass-produce animal-headed humans by the
reel-full, having never remembered anything.

\vspace{4pt}\par
\noindent There is a second debt to the same hobby, and it is a fact about audiences rather than
about iconography. Outside a university lecture theatre, the only places this author has ever been
invited to deliver this material were two summer furry conventions in Alberta, in 2018 and again in
2019. Neither was a sympathetic slot. Both times the ufology hour was programmed directly against
the dance competition --- the single largest event on a furry convention schedule, the thing most of
the building has spent months rehearsing for --- which is, in programming terms, a polite way of
burying a panel. The room filled anyway, and then filled past its chairs: standing along the walls,
sitting on the floor at the front, the doors left open because people were listening from the
corridor. Nobody left early for the music.

\vspace{4pt}\par
\noindent The comparison with the campus club is not flattering to the academy and it is worth
stating without decoration. On a registered university campus, with a chartered club, a room
booking and over a hundred names on a mailing list, the first lecture drew six people. At a
convention of people in animal costumes, opposite the biggest event of the weekend, it was
standing room only. The difference is not intelligence and it is not curiosity; both audiences had
plenty of each. The difference is stigma, and which of the two rooms had already paid it. A furry
convention is by definition a hall full of people who have decided in advance that they do not mind
being seen enjoying something the outside world finds ridiculous. That is precisely the entry fee
this subject charges, and it explains a great deal about who is willing to sit down and examine the
evidence properly. The fandom was not an easier audience. It was a braver one.

\vspace{4pt}\par
\noindent One character type in that hobby taught him more than the rest of it put together, and it
is worth naming because the literature leans on it so heavily. The most-drawn creature in the
entire fandom is not a wolf or a fox. It is the \hi{dragon}. Dragon-kin outnumber every real
animal in the reference folders: horned, winged, plated, feathered, four-legged and two,
western and eastern, in every colour a screen can produce --- tens of thousands of them a year,
none of them drawn from life, because there is nothing alive to draw them from. That is the
control case the jackal-headed argument never had. With a jackal you can at least argue that
somebody once saw a jackal. A dragon has no stock animal behind it at all; it is assembled out
of parts of other animals by people who can itemise exactly which parts they took, and it is
still the single most productive figure in the whole visual culture. A creature can therefore be
drawn more often, more consistently and more lovingly than any real species on Earth, and still
never have existed.

\vspace{4pt}\par
\noindent Which is the reason he cannot read the \textsc{Alpha Draconians} of Section~\ref{sec:racesaz}
with a straight face. Ten feet of horned, winged, armoured reptile from Thuban, offered as the
aristocracy of the reptilian material and as the origin of every dragon myth on Earth: crocodile
hide, rhinoceros horn and the wings of a fruit bat, which is to say the standard mediaeval dragon
kit, assembled the same way the hobby assembles it and from the same parts bin. The claim runs
backwards. The myths are not a memory of the creature; the myths are where the description came
from, and they were being illuminated on European manuscripts centuries before anybody thought
to give the animal a home star. He has spent twenty years watching that exact figure be built
from scratch, in public, by people who make no claim on it whatsoever. It is very difficult to
then be told that the same silhouette, with a star catalogue number attached, is testimony.

\vspace{4pt}\par
\noindent Which brings me, finally, to the debt I left unpaid in Chapter~\ref{ch:life}. I told you
there that one card in the taxonomy of Section~\ref{sec:racesaz} stood apart from the other
thirty-nine, that it was almost certainly not the one you had picked, and that I could not
explain why until you had read this far. Here is the lens, and it is the one I have just
finished describing.

\vspace{4pt}\par
\noindent Run the chain in order. You begin with \hi{animals} --- real ones, the lion and the mantis
and the lungfish and the ant, the actual stock every entry on that list was cut from. Then
comes what a mind does with an animal when it wants to make it a person: it puts the head on
an upright human frame and gives it a name and a voice. That is the \hi{furry} operation, and
I have watched it performed ten thousand times by people who could tell you to the pixel
which parts they invented. Then look at the \hi{aliens} --- the lion-faced Lyran, the mantid,
the fish-skinned Nommo, the Ant People, the Blue Avians in their feathers --- and you are
looking at the identical operation, performed by people who believe the result is a
photograph. And at the far end sit the \hi{alien human--animal hybrids}, the abduction
literature's insistence that these figures cross with us, which is the same operation once
more with the seams showing. Animals, to furries, to aliens, to hybrids: one continuous
procedure, and only the last two steps claim to be memory.

\vspace{4pt}\par
\noindent Thirty-nine of the cards are that procedure. Which is why the exception is the
\textsc{Greys}, and why its \textsc{Earth stock} line is the only blank one in the book. The
Grey is the single figure in the entire taxonomy that was not built by putting an animal's head
on a human body --- because there is no animal in it at all. It is us, caught at the two
moments a human face stops looking human: the foetus, and the face at the far end of ageing.
The skull grown as the body wastes; the eyes enlarged as the features recede; hairless,
sexless, without a mouth worth the name. Every other race on that list is a creature in
costume. The Grey is a mirror, held up at the beginning of a life and again at the end.

\vspace{4pt}\par
\noindent And that, not the zoology, is the most damning entry of the lot. An animal-headed alien only
tells you that somebody had seen an animal. A Grey tells you where the image was
\emph{generated} --- inside, by a human nervous system, with no visitor required and no craft
to get one here. It is precisely the figure you would predict of an experience that arrives at
three in the morning, in a bedroom, to a person who cannot move: not a species, but a face
assembled out of the only face the brain has ever had unlimited access to. Thirty-nine of these
beings were borrowed from the zoo. The fortieth was borrowed from the man in the
mirror, and he is the one I have spent twenty years drawing.

His interests outside software run to Alberta infrastructure --- he is behind the Prairie Lanes
rail and transit concept --- and to the odd corner of provincial history where those interests meet
this one. Alberta is, after all, the home of the world's first UFO landing pad, built at St.\ Paul
in 1967 for Canada's centennial and still standing; the author has publicly proposed that the town
lean into it with a UFO-themed skatepark. He remains open to the possibility that something is
flying around up there. He is simply not willing to accept it on the strength of a blurry
photograph.

\vspace{16pt}
\begin{center}
{\color{ufogreen}\rule{0.34\textwidth}{2.2pt}}\\[13pt]
{\footnotesize\itshape\color{steel} Corrections, better sources and well-argued disagreement
reach the author\\ faster than any of the addresses in Chapter~\ref{ch:media} will reach anyone else.}
\end{center}
